\section*{Chapter \ref{ch:mx:empirical-facts-of-order-books} --- Empirical Facts of Order Books}

\begin{solution}{exo:mx:empirical-facts-of-order-books:1}
Three ticks (the best bid minus the price, 100.00 minus 99.97); zero for a buy at 100.00, which joins the best.
\end{solution}

\begin{solution}{exo:mx:empirical-facts-of-order-books:2}
$0.96+1.64\times0.9=2.44$ ticks.
\end{solution}

\begin{solution}{exo:mx:empirical-facts-of-order-books:3}
\omterm{def:mx:order-types-and-their-uses:hidden}{Hidden orders} are invisible until they trade: a feed shows hidden liquidity only through its executions, so the observable share is that of trades.
A share of orders would measure how much interest is hidden, which only the venue (or its regulator's order-level data) can see.
\end{solution}

\begin{solution}{exo:mx:empirical-facts-of-order-books:4}
About 540. The lifetable estimates the distribution of time to cancellation as if executions did not happen (they censor the clock): it is the
law of an order's patience, not the share of all orders cancelled that quickly; orders that execute early would have been cancelled later in
proportion to the estimated hazard.
\end{solution}

\begin{solution}{exo:mx:empirical-facts-of-order-books:5}
Rising share prices made round lots of 100 shares expensive, so more orders are smaller than a lot; algorithms slice orders finely, and retail
brokers route fractional and small orders. A feed handler or a consolidated view that ignores odd lots misses most trades and, where odd lots rest at
better prices, the true best price.
\end{solution}

\begin{solution}{exo:mx:empirical-facts-of-order-books:6}
When the spread is already one tick and most of the time cannot narrow, liquidity providers compete through queue length and adverse selection shows in
depth and queue dynamics instead of in the spread; the relation holds for the implicit spread of chapter~6, not for the quoted one.
\end{solution}

\begin{solution}{exo:mx:empirical-facts-of-order-books:7}
The simulated day: 1.8\%, 4.6\% and 12.4\% executed within 1, 10 and 60 seconds, against 3.4\%, 8.7\% and 15.0\% for large US stocks. Executions in the
simulator are about half as fast at a second and close at a minute, while its cancellations are a hundred times slower: the simulated market is slow at
revising quotes, not at trading, so it overstates how long a resting order survives the arrival of news.
\end{solution}

\begin{solution}{exo:mx:empirical-facts-of-order-books:8}
The simulator's orders live a hundred times longer than real ones and its liquidity providers cancel slowly: a resting order there faces much less
competition to cancel stale quotes, and faster traders who pick off slow quotes are absent. Fill rates and adverse selection must be measured against a
market whose cancellation speed matches the venue's (chapter~27), or live at small size.
\end{solution}

\begin{solution}{pb:mx:empirical-facts-of-order-books:1}
\textbf{1.} One tick behind the best (1\,508 against 1\,366 shares in chapter~1's hour); real books are humped too (Bouchaud, Mézard and Potters).
\textbf{2.} 0.090 and 0.044 per share per second: stale orders at the best are pulled when the efficient price moves through them.
\textbf{3.} 0.5\%, 18.6\% and 40.1\%.
\textbf{4.} A power law with a cumulative exponent near 1.5 up to about 2\,000 ticks (London); the simulator places orders on ten levels with geometrically
falling rates, none beyond nine ticks.
\textbf{5.} 96.6\% of the time; mean 1.042 ticks.
\textbf{6.} $s\approx0.96+1.64\,\sigma_1$, $R^2=0.989$.
\textbf{7.} $R^2=0.48$, slope 0.19: a large-tick book whose spread cannot move.
\textbf{8.} U-shaped: 92\,200 shares in the first thirteenth, 24\,700 at the lowest, 55\,600 in the last; eight sessions averaged.
\textbf{9.} Durations from the add to each cancellation or trade, restarting after partial events, estimated by the lifetable method with executions censoring
cancellations.
\textbf{10.} Large stocks 54\% within 0.1~s and 74\% within 1~s; small stocks 37\% and 49\%.
\textbf{11.} 0.8\%, 6.8\% and 40.0\%.
\textbf{12.} 11.5\%.
\textbf{13.} Median 0.19 seconds back to a one-tick spread; 92.4\% within a second (750 events).
\textbf{14.} 15.5 against 16.6 in 2026: close.
\textbf{15.} \emph{Named result}: the simulator reproduces the humped depth profile, the U-shaped intraday volume (planted), fast resilience and the
cancel-to-trade ratio; it fails the power law of \omterm{def:mx:empirical-facts-of-order-books:relative}{relative limit prices}, the spread--volatility relation (large tick), the opening spread, cancellation
speed (a hundred times too slow) and hidden liquidity (none).
\textbf{16.} Slow cancellations and the absence of fast competitors make resting orders look safer and fill more benignly: fill rates biased up, adverse
selection down.
\textbf{17.} The thin book beyond nine ticks: large orders would walk into nothing, overstating impact far from the best.
\textbf{18.} Cancellations per trade fell from 19.8 to 16.6; the hidden rate rose from 10.5\% to 32.2\%; the odd-lot rate from 19.2\% to 68.4\%.
\textbf{19.} Cancellation speed (fast quote revisions after efficient-price moves), then \omterm{def:mx:order-types-and-their-uses:hidden}{hidden orders} and deeper placement.
\textbf{20.} A test any model of the market must pass before its answers are trusted.
\end{solution}

\begin{solution}{iq:mx:empirical-facts-of-order-books:1}
The best level is where executions and stale-quote cancellations happen first; behind it, orders are out of reach of most \omterm{def:mx:the-limit-order-book:orders}{market orders} and accumulate,
refreshed by new arrivals. \iqlookfor{flows by distance: execution, cancellation and arrival rates.}
\end{solution}

\begin{solution}{iq:mx:empirical-facts-of-order-books:2}
That adverse selection sets it: if liquidity provision earns at most a marginal profit, the spread must cover the price move each trade causes, and that
move is the \omterm{def:mx:empirical-facts-of-order-books:vol}{volatility per trade}. \iqlookfor{Wyart et al.'s break-even argument and the role of trade impact in volatility.}
\end{solution}

\begin{solution}{iq:mx:empirical-facts-of-order-books:3}
Survival analysis: treat execution as censoring for the cancellation clock (and vice versa), and estimate with Kaplan--Meier or a lifetable; restart the
clock after partial events if the question is about each decision. \iqlookfor{recognising censoring; a naive mean over cancelled orders is biased.}
\end{solution}

\begin{solution}{iq:mx:empirical-facts-of-order-books:4}
List the facts the strategy's P\&L depends on (queue dynamics, cancellation speed, adverse selection, resilience), measure each on the simulator and on
real data or published statistics, and backtest only once they agree within tolerances; test that the conclusion survives the facts the simulator gets
wrong. \iqlookfor{a validation plan tied to the strategy, not a generic realism score.}
\end{solution}

\begin{solution}{iq:mx:empirical-facts-of-order-books:5}
Uncertainty is highest after the overnight news: liquidity providers face more informed flow and more price risk, and the book is still filling after the
opening auction; as information is absorbed, competition narrows the spread. \iqlookfor{adverse selection and inventory risk varying through the day.}
\end{solution}

\begin{solution}{iq:mx:empirical-facts-of-order-books:6}
Its sender is searching (for hidden liquidity or a counterparty) or tracking a moving price, rather than offering liquidity. Test: condition later price
moves and executions on fleeting-order activity and compare with long-lived orders; check whether \omterm{def:mx:empirical-facts-of-order-books:life}{fleeting orders} cluster before trades on the same side.
\iqlookfor{a testable hypothesis with a control group, not a story.}
\end{solution}
